\section{The Evidence Surface}
\label{sec:framework-15}

The surface partition of Section~\ref{sec:framework-4} is exhaustive, and
Theorem~1 proves it. This section argues that the theorem is true and the axiom it
rests on is incomplete, and that closing the gap requires a sixth surface rather than
a wider reading of an existing one.

The argument is put here, at the end of the chapter, because it is a critique of the
foregoing rather than a qualification of it. Nothing above needs to be withdrawn.

\subsection{What the theorem actually proves}

Theorem~1 establishes that every element of every interaction is assigned to exactly
one control surface. Its first proof step is the load-bearing one: \emph{by Axiom 1,
every interaction is described by} $(x, u, M, \theta, o)$. The partition is then
exhaustive over that tuple, and the conclusion follows.

So the theorem's scope is the tuple, and its truth is relative to the tuple. It proves
completeness of the partition, not completeness of the tuple. Axiom~1 is careful about
this and says so directly: it claims governance scope rather than completeness, and
concedes that elements outside the tuple ``may influence outputs'' while declining to
model them.

Among the elements named as outside: \emph{retrieval index state}. The three-tier
classification places it in Tier~2, known but unmodeled, contributing documented
ungoverned risk.

\subsection{Why that placement does not hold for retrieval-augmented systems}

Tier~2 is the right classification for infrastructure that perturbs an output without
being part of what the output claims. Load balancing, caching, replica differences: an
interaction is not \emph{about} its own routing, and nothing an auditor asks concerns
which replica served the request.

Retrieved evidence is not of that kind. When a system answers a regulated question from
a retrieved instrument, the instrument is not a factor that influenced the output. It is
the \emph{basis of the assertion} --- the thing an auditor is asking about when they ask
why. Placing it in Tier~2 makes the governance framework unable to represent the one
element the audit is for.

The chapter's own architecture shows the consequence. The prompt surface governs the
assembled artifact, and the evaluation engine operates on that assembled text rather
than on the materials that composed it; the composition process is invisible to the
evaluation layer. So retrieved evidence enters the governed tuple as opaque tokens
inside $x$. Two prompts of identical text, one assembled from a current authoritative
instrument and one from a superseded foreign-jurisdiction instrument, are the same
governed artifact under this model. They evaluate identically, produce identical audit
records, and are distinguishable only by information the framework has declared out of
scope.

That is not a gap in enforcement. It is a gap in \emph{representation}, and it cannot be
closed by strengthening any of the five surfaces, because none of them has a term for a
source.

\subsection{The sixth surface}

Adding evidence therefore requires amending Axiom~1 rather than extending
Definition~3.2. Write the amended tuple
\begin{equation}
  i = (x, u, M, \theta, o, E)
  \label{eq:meta-amended-tuple}
\end{equation}
where $E$ is the evidential basis available to the interaction: the identified,
versioned sources the assertion is permitted to rest on. The partition then gains
\begin{equation}
  \sigma_{\mathrm{evidence}} = \{E\}
  \label{eq:meta-evidence-surface}
\end{equation}
and Theorem~1 is restated over the six-element tuple, where it holds by the same proof.

Its governability class is worth stating precisely, because it is the reason the surface
is worth adding:
\begin{equation}
  \gov(\sigma_{\mathrm{evidence}}) =
  \begin{cases}
    \text{full}    & \text{retrieval-provided evidence, identified at inference time}\\
    \text{external}& \text{parametric evidence absorbed in training}
  \end{cases}
  \label{eq:meta-evidence-governability}
\end{equation}

This is the only surface in the framework whose class is \emph{split by deployment
choice rather than fixed by the technology}. The prompt surface is governable because
prompts are constructed; model internals are external because they are not observable.
Evidence is either, depending on whether the architecture makes it explicit --- which
turns a governability classification into a design requirement, and is the strongest
architectural conclusion in this book.

\subsection{What the new surface requires}

Three things the existing five do not supply.

\paragraph{An applicability relation.} The output-evaluation machinery of the runtime
control layer is defined over $\mathrm{applicable}(R_{\mathrm{ROC}}, i)$, and the
completeness theorem for that layer asserts its audit record is total over exactly that
set. The function is used and never defined anywhere in this framework. For constraints
the omission is survivable, because applicability of a rule to an interaction is
settled by level and domain. For evidence it is not: whether a source governs a case is
the substance of the question, and it is developed in
Chapter~\ref{ch:grounding-context} as a relation to a resolved context of jurisdiction,
valid time, material facts, and intended use.

\paragraph{A second time axis.} Axiom~5 gives temporal identity over a single totally
ordered domain, which is sufficient for asking which rule version was in force when an
interaction ran. Evidence needs two axes: the time a source was valid, and the time the
system could have known it. Their divergence is a detectable failure in both
directions, and neither is expressible with one axis.

\paragraph{A dependence condition.} The prompt surface can establish that required text
was present and the output surface that prohibited text was absent. Neither can establish
that an assertion \emph{rested on} the evidence cited for it. The framework's own worked
example makes the point: a citation check fails, the prompt is reinforced, the model
returns the same analysis with citations attached, and the second evaluation passes.
Nothing in the five surfaces can distinguish that from grounded assertion, because a
presence check on the artifact is the only instrument available.

\subsection{Consequence for the guarantees}

The composed guarantee of Section~\ref{sec:framework-12} is unaffected in what it claims and
narrowed in what it covers. Every theorem in this chapter remains true over the tuple it
quantifies over. What changes is the residual-risk accounting: evidence moves from Tier~2
undocumented influence to a governed surface with its own class, its own audit fields, and
its own failure modes --- which means the risk attributable to it becomes nameable
instead of being absorbed into the irreducible remainder.

The conditions on $\sigma_{\mathrm{evidence}}$ are the grounding conditions of
Chapter~\ref{ch:taxonomy}; the invariants that make them checkable, and the
measurement that establishes them, are stated where the other system invariants are. This section
claims only that the surface belongs in the partition, and that its absence was a
scoping decision made before retrieval-augmented deployment was the normal case.
